\section{The finite-trace embedding theorem}
\label{sec:master}

\subsection{Trace systems}

All run-time parameters---the original input, requested precision, finite
bounds, and externally supplied certificates---may be paired into one natural
number \(u\).  A trace system consists of fixed primitive-recursive functions
\[
\begin{array}{ll}
 \mathsf{init}(u), & \mathsf{bound}(u),\\
 \mathsf{step}(u,i,s), & \mathsf{inv}(u,i,s)\in\{0,1\},\\
 \mathsf{good}(u)\in\{0,1\}, & \mathsf{accept}(u,s)\in\{0,1\}.
\end{array}
\]
Define the run by primitive recursion:
\[
 \mathsf{run}(u,0)=\mathsf{init}(u),\qquad
 \mathsf{run}(u,i+1)=\mathsf{step}(u,i,\mathsf{run}(u,i)).
\]
The intended proof obligations are the following quantifier-free formulas:
\begin{align}
 \mathsf{good}(u)=1&\to
   \mathsf{inv}(u,0,\mathsf{init}(u))=1, \tag{T0}\label{eq:T0}\\
 \mathsf{good}(u)=1\wedge \mathsf{inv}(u,i,s)=1&\to
   \mathsf{inv}(u,i+1,\mathsf{step}(u,i,s))=1, \tag{T1}\label{eq:T1}\\
 \mathsf{good}(u)=1\wedge
 \mathsf{inv}(u,\mathsf{bound}(u),s)=1&\to
   \mathsf{accept}(u,s)=1. \tag{T2}\label{eq:T2}
\end{align}

\begin{theorem}[Finite-trace reflection]
\label{thm:trace-reflection}
If \(\ipra\) proves \eqref{eq:T0}--\eqref{eq:T2}, then it proves
\[
 \forall u\;\bigl(\mathsf{good}(u)=1\to
 \mathsf{accept}(u,\mathsf{run}(u,\mathsf{bound}(u)))=1\bigr).
\]
\end{theorem}

\begin{proof}
Keep \(u\) free.  Apply quantifier-free induction to
\[
 A(i)\;:\equiv\;
 \mathsf{good}(u)=1\to
 \mathsf{inv}(u,i,\mathsf{run}(u,i))=1.
\]
This is quantifier free: all functions are fixed primitive-recursive symbols,
and implication is a propositional connective rather than a quantifier.  The
base case is \eqref{eq:T0} after rewriting \(\mathsf{run}(u,0)\).  For the
successor case, assume \(A(i)\) and \(\mathsf{good}(u)=1\).  The induction
hypothesis gives the invariant at \(\mathsf{run}(u,i)\); \eqref{eq:T1} gives
the invariant after one step; the defining equation for \(\mathsf{run}\)
identifies that state with \(\mathsf{run}(u,i+1)\).

Specialize the resulting universal statement to
\(i=\mathsf{bound}(u)\).  Apply \eqref{eq:T2} with
\(s=\mathsf{run}(u,\mathsf{bound}(u))\).  Finally introduce \(u\)
universally.  No induction formula contains an unbounded quantifier.
\end{proof}

\begin{remark}
The theorem is intentionally modest.  It does not say that every true
primitive-recursive assertion is provable.  It says that a supplied base case,
local invariant proof, and final implication compose into the expected bounded
algorithm theorem.  The mathematical work lies in choosing an invariant for
which \eqref{eq:T1} is actually derivable.
\end{remark}

\subsection{Finite query tables}

A represented input may have a global contract such as
\[
 \mathsf{Reg}_a(p)\;:\equiv\;
 \forall m,n\;R(a(p,m),a(p,n),m,n).
\]
Putting \(\mathsf{Reg}_a(p)\) inside the induction predicate would violate the
intended proof discipline.  Instead, the trace declares a finite query list
\(q(u,0),\ldots,q(u,L(u)-1)\).  Define \(\mathsf{GoodTable}(p,u)=1\) as the
Boolean conjunction of exactly the regularity, domain, and magnitude facts
needed about those samples.

\begin{lemma}[Finite instantiation]
\label{lem:finite-instantiation}
For every fixed query generator, \(\ipra\) proves
\[
 \mathsf{Reg}_a(p)\to\mathsf{GoodTable}(p,u)=1,
\]
provided each field of \(\mathsf{GoodTable}\) is an instance or a finite
algebraic consequence of the global contract.
\end{lemma}

\begin{proof}
The query list is finite and primitive recursive.  For each entry, eliminate
the universal quantifiers in \(\mathsf{Reg}_a(p)\) at the displayed numerical
indices.  The remaining rational equalities and inequalities are combined by
the Boolean-fold correctness lemma.  If later queries depend on earlier
answers, perform the same argument along the actual trace: at stage \(i\), the
query is a primitive-recursive term in the previous state, and universal
elimination permits that term.  The accumulator invariant says that all facts
queried so far have been inserted into the table.  This invariant is a Boolean
equation and is therefore quantifier free.
\end{proof}

This lemma explains how universally quantified name hypotheses can remain in
the surrounding context while the induction formula itself stays finite.

\subsection{Fix precision, then quantify}

\begin{theorem}[Fixed-precision compilation]
\label{thm:fixed-precision}
Let \(k\) be a requested output precision and let \(u=\angles{p,k,c}\) encode
all numerical inputs and finite certificates.  Suppose:
\begin{enumerate}
\item global hypotheses \(H(p)\) imply \(\mathsf{good}(u)=1\) by finite
      instantiation;
\item the trace system satisfies \eqref{eq:T0}--\eqref{eq:T2};
\item acceptance implies a finite output observation
      \(\mathsf{Out}(p,k,A(p,k))\).
\end{enumerate}
Then \(\ipra\) proves
\[
 H(p)\to\forall k\;\mathsf{Out}(p,k,A(p,k)).
\]
\end{theorem}

\begin{proof}
Keep both \(p\) and \(k\) free.  From \(H(p)\), derive the finite good table
for this particular \(k\).  Apply \cref{thm:trace-reflection}, then the final
soundness implication.  At no point is induction applied to a formula
containing \(\forall k\).  Once the finite derivation has been obtained with
\(k\) free, apply universal introduction to \(k\).
\end{proof}

This simple reordering resolves a large class of reported ``iPRA induction is
too weak'' concerns.  Frequently the problem is not the estimate but the
attempt to induct on its completed all-precision interpretation.

\subsection{Certificate reflection}

A computation need not be rerun inside the checker.  Let a certificate
\(c\) contain a state list \(s_0,\ldots,s_B\).  Define
\(\mathsf{TraceOK}(u,c)=1\) to assert
\[
 s_0=\mathsf{init}(u),\qquad
 s_{i+1}=\mathsf{step}(u,i,s_i)\ (i<B),\qquad
 \mathsf{inv}(u,i,s_i)=1\ (i\le B),
\]
and the final acceptance equation.  Every clause is checked by a finite fold.

\begin{corollary}[Uniform finite-certificate checker]
\label{cor:certificate-checker}
For a fixed trace language, \(\ipra\) proves
\[
 \mathsf{TraceOK}(u,c)=1\to\mathsf{Conclusion}(u,\mathsf{output}(c)).
\]
\end{corollary}

\begin{proof}
Use fold elimination to obtain every local transition and invariant clause.
The final clause directly supplies acceptance, or alternatively reconstruct
the run equality by induction on the list index and invoke
\cref{thm:trace-reflection}.  Both routes use only finite list induction with a
Boolean invariant.
\end{proof}

Certificate reflection is the right generic mechanism for externally produced
polynomial normalizations, isolating regions, factorization traces, finite
covers, and large numerical constant tables.

\subsection{Composition theorem}

\begin{theorem}[Composition of certified transducers]
\label{thm:composition}
The sequential composition of two \(\CFP\) transducers is a \(\CFP\)
transducer.  Its precision transformer, trace bound, and checker are primitive
recursive, and its soundness theorem is provable in \(\ipra\) from the two
component soundness theorems.
\end{theorem}

\begin{proof}
Encode a composite state as a phase bit together with the current component
state and the finite metadata passed between components.  In phase zero, run
the first transducer for its primitive-recursive bound.  Its final state
contains the output code and every certified magnitude/domain bound needed by
the second precision transformer.  Switch the phase bit, initialize the
second state, and run it for its bound.  The total bound is the sum of the two
bounds plus one switch step.

The composite invariant is a Boolean case split on the phase bit.  In phase
zero it is the first invariant.  At the switch it includes the first
acceptance predicate and the second input-good predicate.  In phase one it is
the second invariant together with the frozen metadata.  Each case is a
quantifier-free equation.  Apply finite-trace reflection.  The final soundness
implication is exactly the second component theorem after substituting the
certified first output.  Primitive-recursive closure under pairing,
composition, case distinction, addition, and primitive recursion supplies all
new symbols.
\end{proof}

The same construction handles finite acyclic proof networks and bounded
adaptive calls.  A network is flattened into one phase-indexed trace.

\subsection{Export to the native calculus}

Once the endpoint
\(\Gamma\vdash[\texttt{PRAAx}]\varphi\) has been constructed, the existing
repository theorem gives
\[
 \mathsf{NPrf}(\Gamma\text{ translated},\varphi\text{ translated}).
\]
No additional mathematical proof is required at this stage.  Conversely, the
native public/guarded calculi can serve as compact certificate targets when
their formula restrictions are convenient.  What must be avoided is invoking
bi-interpretability before there is any object-language proof to transport.
